\documentclass[../main.tex]{subfiles}

\graphicspath{{../image/}{../figure/}}

\begin{document}

\chapter{Cauchy 问题}\label{chapter:ch04-cauchy问题}

双曲守恒律 Cauchy 问题的理论面临两大挑战。第一，从光滑初值出发的古典解会自发产生间断；因而一般而言，整体上只存在弱解。第二，Cauchy 问题的弱解不唯一。无须深究即可遭遇这些困难。正如在\cref{section:ch04-经典解的崩溃}、\cref{section:ch04-弱解的非唯一性}与\cref{section:ch04-Euler方程}中所示，即使在单个或多个空间维数中最简单的非线性双曲守恒律层面，它们也会出现。

弱解的 Cauchy 问题将在\cref{section:ch04-弱解}中表述。为克服非唯一性的障碍，须施加限制以剔除不稳定的、物理上无关的或其他不合需要的解，期望以此挑选出唯一的可容解。本章将引入两种可容性判据：要求可容解满足指定熵不等式；以及要求可容解为含扩散项方程组解族在扩散渐近消失时的极限。我们将对这两种判据作初步比较。

关于识别数学上与物理上有意义的边界条件这一问题的初步讨论，将在\cref{section:ch04-初边值问题}中给出。

本章最后一节\cref{section:ch04-Euler方程}收集三维空间中（等熵）气体动力学 Euler 方程组的一组代表性结果，它们来自较早或新近的工作，并展示双曲守恒律这一重要领域中当前的研究趋势。

\subfile{Chapter-04/section-01.tex}
\subfile{Chapter-04/section-02.tex}
\subfile{Chapter-04/section-03.tex}
\subfile{Chapter-04/section-04.tex}
\subfile{Chapter-04/section-05.tex}
\subfile{Chapter-04/section-06.tex}
\subfile{Chapter-04/section-07.tex}
\subfile{Chapter-04/section-08.tex}
\subfile{Chapter-04/section-09.tex}

\end{document}
